% 第1章 1.2.3：积分正交由乘积的整体积分决定。[-1,1] 上 Lebesgue 测度。
% (a) 1 与 x：乘积 x 的正负面积各 1/2，相加为 0，二者正交。
% (b) 1 与 x^2：乘积 x^2 的积分为 2/3，不正交。
\begin{tikzpicture}[font=\figurefont\footnotesize]
\pgfplotsset{every axis/.style={width=46mm,height=36mm,axis lines=middle,
  x axis line style={draw=FigureStroke,line width=.6pt,-{Stealth[length=2mm]}},
  y axis line style={draw=FigureStroke,line width=.6pt,-{Stealth[length=2mm]}},
  tick style={draw=FigureStroke},tick label style={font=\figurefont\scriptsize,text=FigureInk},
  xmin=-1.15,xmax=1.25,ymin=-1.15,ymax=1.3,xtick={-1,1},ytick=\empty,
  xlabel={$x$},xlabel style={anchor=north west,font=\figurefont\scriptsize},
  title style={font=\figurefont\footnotesize},samples=100,axis on top}}
\begin{axis}[name=A,title={(a) $1\cdot x$，积分 $0$}]
\addplot[fill=FigureBlueFill,draw=none,domain=0:1] {x} \closedcycle;
\addplot[fill=FigureGreenFill,draw=none,domain=-1:0] {x} \closedcycle;
\addplot[draw=FigureBlue,line width=1.1pt,domain=-1:1] {x};
\node[font=\figurefont\scriptsize,text=FigureBlue] at (axis cs:0.65,0.28) {$+\tfrac12$};
\node[font=\figurefont\scriptsize,text=FigureGreen] at (axis cs:-0.65,-0.28) {$-\tfrac12$};
\end{axis}
\begin{axis}[at={(A.east)},anchor=west,xshift=12mm,title={(b) $1\cdot x^2$，积分 $\tfrac23$}]
\addplot[fill=FigureBlueFill,draw=none,domain=-1:1] {x^2} \closedcycle;
\addplot[draw=FigureBlue,line width=1.1pt,domain=-1:1] {x^2};
\node[font=\figurefont\scriptsize,text=FigureBlue] at (axis cs:0.72,0.22) {$+\tfrac13$};
\node[font=\figurefont\scriptsize,text=FigureBlue] at (axis cs:-0.72,0.22) {$+\tfrac13$};
\end{axis}
\end{tikzpicture}
